\uuid{ZHM4}
\titre{Dérivée sur un intervalle}
\auteur{bodin ; hanani ; mzari}
\organisation{exo7}
\licence{CC BY-NC-SA 4.0}
\datecreate{2020-01-20}
\niveau{L1}
\module{Analyse}
\chapitre{Dérivabilité des fonctions réelles}
\sousChapitre{Autre}
\difficulte{3}
\theme{Dérivées, dérivée sur un intervalle}
\exo7qcmid{lille-1.derivees.021}

\contenu{
\begin{question}
\qclassification{124.00}
\qtags{motcle=dérivée sur un intervalle}
Soit $f(x)=\arctan x+\arctan \frac{1}{x}$ définie sur $\Rr^*$. Quelles sont les bonnes réponses ?
\begin{answers}
    \good{$\forall x\in \Rr^*$, $f'(x)=0$}
    \bad{$\forall x\in \Rr^*$, $\displaystyle f(x)=\frac{\pi}{2}$}
    \bad{La fonction $f$ est paire.}
    \good{$\displaystyle f(x)=\frac{\pi}{2}$ si $x>0$ et $\displaystyle f(x)=-\frac{\pi}{2}$ si $x<0$}
\end{answers}
\begin{explanations}
La fonction $f$, tout comme la fonction $\arctan$, est impaire. On calcule $f'(x)$ pour $x\neq 0$ :
$$f'(x)=\frac{1}{1+x^2}+\frac{\left(\frac{1}{x}\right)'}{1+\left(\frac{1}{x}\right)^2}=0$$
Donc $f$ est constante sur chaque intervalle de son domaine de définition :
$$f(x)=\left\{\begin{array}{ll}\displaystyle f(1)=\frac{\pi}{2}&\mbox{si }x>0\\ \\ f(-1)=-\frac{\pi}{2}&\mbox{si }x<0.\end{array}
\right.$$
\end{explanations}
\end{question}
}
